% ============================================================
% Abstract
% ============================================================
\chapter*{Abstract}
\addcontentsline{toc}{chapter}{Abstract}

Learners' affective and cognitive states---engagement, confusion,
frustration, and boredom---shape learning outcomes, yet in online settings
they are largely invisible to instructors. This capstone designed, built, and
deployed a browser-based intelligent tutoring platform that makes those
states measurable while students learn with retrieval-grounded,
large-language-model (LLM) tutoring. The platform combines (i) a
course-authoring and RAG (retrieval-augmented generation) pipeline with
hybrid dense--sparse retrieval; (ii) two learner-facing modes---structured
lessons with a docked tutoring chatbot, and a free-form dialogue workspace
over learner-uploaded sources; (iii) four theory-grounded, learner-initiated
strategy interventions (practice testing, interrogative elaboration,
stepwise scaffolding, and SM-2 spaced repetition); and (iv) a multimodal
sensing layer that time-aligns webcam-based facial analysis (eight-class
emotion probabilities from OpenFace~3.0 and eighteen facial action-unit
intensities from Py-Feat), WebGazer eye gaze, an exploratory pupil-size
proxy, and dense interaction telemetry on a single wall-clock anchor.
An area-of-interest (AOI) attention layer derived from the SEEV tradition
scores how gaze allocation aligns with task-expected attention, and a
session-replay environment with retrospective qualitative coding turns each
learning session into an inspectable multimodal record.

Sensing is deliberately decoupled from pedagogy: interventions are
exclusively learner-initiated and are never triggered by affective signals,
so affect measurements form a clean observational layer rather than a
confound in the learning loop. The measurement methodology comprises
cross-modal convergent validation of facial-expression-derived engagement
estimates against gaze-on-AOI dwell as a behavioral reference criterion;
a nine-construct LLM text-mining pipeline for executive-function and
affect markers in learner utterances; a mandatory five-option emotion
self-report probe administered every fifteen minutes; and a parallel-form
pre/post assessment design for learning gain.

A user study with \todo{N} participants evaluated the platform.
\todo{One-sentence summary of the validation result (agreement between
system-detected and self-reported states).} \todo{One-sentence summary of
the learning-gain result with effect size.} The report documents the
system's architecture and data model in full, states its limitations
candidly---including known data-quality caveats and the boundary between
validated and exploratory signals---and argues that the platform's principal
contribution is a working, privacy-conscious instrument for studying
affect, attention, and strategy use in authentic online learning.

\vspace{1em}
\noindent\textbf{Keywords:} intelligent tutoring systems; affective
computing; multimodal learning analytics; retrieval-augmented generation;
eye tracking; facial action units; learning strategies; spaced repetition.
\clearpage
